\documentclass[11pt,a4paper]{article}
\usepackage[utf8]{inputenc}
\usepackage{amsmath}
\usepackage{lmodern}

\title{Lorem Ipsum as a Placeholder Corpus for Scientific Typesetting}
\author{Martin Petkovski}
\date{July 2026}

\begin{document}
\maketitle

\begin{abstract}
Lorem ipsum is commonly treated as meaningless filler, yet its stable rhythm makes it a useful test corpus for scientific publishing systems. This draft evaluates whether a static web renderer can preserve the structure of an Overleaf-compatible \LaTeX document while presenting it with conventional academic typography.
\end{abstract}

\section{Introduction}

Scientific manuscripts combine prose, mathematical notation, hierarchical sections, lists, and references. A useful test document should exercise each of these elements without allowing the subject matter to distract from the typesetting. Lorem ipsum therefore provides a convenient neutral corpus \cite{cicero45}.

The objective of this draft is not to establish a new result. Instead, it tests three practical requirements:

\begin{enumerate}
\item the source remains valid and editable as a conventional \LaTeX file;
\item the public website renders the source as a readable static paper; and
\item the same page produces a clean A4 document when printed.
\end{enumerate}

\section{Method}

Let $w_i$ denote the number of characters in the $i$th placeholder sentence. We define the normalized visual density of a paragraph containing $n$ sentences as

\begin{equation}
D = \frac{1}{n}\sum_{i=1}^{n}\frac{w_i}{W},
\label{eq:density}
\end{equation}

where $W$ is the available line width measured in typographic units. This deliberately simple quantity is sufficient to test inline variables, summation notation, fractions, subscripts, and displayed equations.

\subsection{Test corpus}

The primary sample is the familiar sequence: \emph{Lorem ipsum dolor sit amet, consectetur adipiscing elit}. Additional paragraphs vary sentence length so that justification, indentation, and page flow can be inspected under realistic conditions.

\begin{itemize}
\item Short clauses reveal excessive word spacing.
\item Long clauses reveal poor line breaking and narrow margins.
\item Mixed mathematical and textual content reveals baseline inconsistencies.
\end{itemize}

\section{Observations}

Lorem ipsum dolor sit amet, consectetur adipiscing elit. Integer finibus, augue id posuere volutpat, sem justo fermentum libero, vitae faucibus arcu nibh sed neque. Nulla facilisi. Sed dignissim lectus non sapien feugiat, vitae malesuada metus faucibus.

Praesent consequat, nibh sed varius tincidunt, justo velit tincidunt lorem, vel volutpat lacus arcu id massa. The layout remains legible when a citation is inserted near the end of a line \cite{knuth1984}. Curabitur consequat lectus sit amet lacus volutpat, quis tincidunt est fermentum.

\subsection{Expected rendering characteristics}

The title should be prominent but restrained, the abstract should use a slightly smaller measure, and body text should remain close to the conventional eleven-point article style. Section headings should establish hierarchy without adopting the visual language of the surrounding website.

\section{Conclusion}

This draft provides a compact regression document for the Draft Papers section. It is intentionally disposable: future papers can replace the placeholder prose while retaining the same Overleaf-compatible workflow, static URL structure, mathematical rendering, and print behavior.

\begin{thebibliography}{9}
\bibitem{cicero45}
M. Tullius Cicero.
\emph{De Finibus Bonorum et Malorum}.
45 BCE.

\bibitem{knuth1984}
Donald E. Knuth.
\emph{The TeXbook}.
Addison--Wesley, 1984.
\end{thebibliography}

\end{document}
